%============================================================
% Chapter 3: Binary Search on Partly Sorted Arrays
% STATUS: PLANNED  (generated by tools/outline.py; replace this file when the chapter is written)
%============================================================
\chapter{Binary Search on Partly Sorted Arrays}
\label{ch:sorted_arrays}

\begin{tcolorbox}[enhanced, colback=white, colframe=secondaryblue, boxrule=1pt, arc=2mm]
\textbf{\large\textcolor{primaryblue}{The array is not sorted, but each half tells you which way to go: rotated arrays, peaks, bitonic arrays and paired structures.}}
\end{tcolorbox}

\chapterfacts{1100}{1900}{3}{6}{30}

\begin{plannedbox}[Planned Contents --- this chapter is not written yet]
\textbf{Part A --- Lessons}
\begin{enumerate}[leftmargin=*, label=\textbf{A\arabic*.}, itemsep=1pt]
    \item Rotated sorted arrays: one half is always sorted; finding the minimum and the rotation point
    \item Duplicates in rotated arrays and the $\bigO{n}$ worst case
    \item Peaks: why a local maximum always exists and how the slope tells you where
    \item Bitonic (mountain) arrays: find the peak, then search both sides
    \item Structure hidden in indices: pairs, parity of positions and missing values
\end{enumerate}
\textbf{Part B --- Patterns} (with the two solved problems of each pattern)
\begin{enumerate}[leftmargin=*, label=\textbf{P\arabic*.}, itemsep=1pt]
    \item \textbf{Rotated arrays} --- Find Minimum in Rotated Sorted Array {\small\color{darkgray}(LeetCode 153)}, Search in Rotated Sorted Array II {\small\color{darkgray}(LeetCode 81)}
    \item \textbf{Peaks and mountains} --- Find Peak Element {\small\color{darkgray}(LeetCode 162)}, Find in Mountain Array {\small\color{darkgray}(LeetCode 1095)}
    \item \textbf{Structure hidden in indices} --- Single Element in a Sorted Array {\small\color{darkgray}(LeetCode 540)}, Kth Missing Positive Number {\small\color{darkgray}(LeetCode 1539)}
\end{enumerate}
\textbf{Connections.} Search in Rotated Sorted Array (LC 33) is solved in the recursion book.
\end{plannedbox}
